\begin{frame}[standout]

    Focus sur Redis

    \note{
        Redis est un exemple concret de base clé-valeur très rapide,
        particulièrement connue pour son fonctionnement en mémoire et ses usages
        de cache, de session et de classement.
    }

\end{frame}

\begin{frame}{Redis : expiration, éviction et persistance}
    \begin{itemize}
        \item Structures courantes : chaînes, listes, ensembles, ensembles triés et hachages ; Redis Cluster distribue les clés en slots de hachage.
        \item Les clés peuvent recevoir un \textbf{TTL} pour expirer automatiquement.
        \item L'\textbf{éviction} libère de la mémoire lorsqu'une limite est atteinte ; elle dépend d'une politique configurée. Expiration et éviction sont distinctes.
        \item \textbf{RDB} enregistre des instantanés périodiques : sauvegarde compacte, mais les écritures récentes peuvent être perdues.
        \item \textbf{AOF} journalise les opérations : reprise plus fidèle selon la configuration, au prix d'un journal plus volumineux et d'une relecture.
    \end{itemize}

    \note{
        Redis fonctionne principalement en mémoire et propose plusieurs
        structures (chaînes, listes, ensembles, ensembles triés et hachages).
        Les clés sont distribuées entre les nœuds de Redis Cluster par slots
        de hachage. Les garanties de reprise dépendent des réglages de persistance
        et de réplication : un cache et un état métier critique n'ont pas les
        mêmes exigences.
    }
\end{frame}

\begin{frame}{Redis : une architecture master-slave}
    \begin{itemize}
        \item Le \textbf{master} reçoit les écritures et les propage à ses \textbf{réplicas} (slaves), de manière asynchrone par défaut ; ceux-ci peuvent servir des lectures.
        \item \textbf{Redis Sentinel} surveille les instances et promeut un réplica en cas de panne du master.
        \item \textbf{Redis Cluster} répartit les slots de hachage entre plusieurs masters, chacun ayant ses propres réplicas.
    \end{itemize}

    \note{
        Dans Redis Cluster, les nœuds échangent leur état par gossip, mais ce
        n'est pas une architecture peer-to-peer comme celle de Cassandra : pour
        un slot donné, seul le master accepte les écritures. La réplication étant
        asynchrone, une écriture acquittée peut être perdue si le master tombe
        avant de l'avoir transmise.
    }
\end{frame}

\begin{frame}{Redis : usages typiques}
    \begin{itemize}
        \item \textbf{Cache} : stocker en mémoire les résultats fréquemment demandés.
        \item \textbf{Sessions utilisateur} : retrouver une session par son identifiant, avec une expiration automatique grâce au TTL.
        \item \textbf{Classements} : maintenir un classement en temps réel avec un ensemble trié, malgré des mises à jour de score très fréquentes.
    \end{itemize}

    \begin{examplebox}{Limites}
        Les données en mémoire coûtent plus cher, nécessitent un mécanisme de persistance ou de reconstruction, et ne se prêtent pas aux requêtes complexes sur d'autres attributs.
    \end{examplebox}
\end{frame}

\begin{frame}{Étude de cas}

    https://redis.io/customers/

    https://redis.io/customers/france-travail/

    Problématique.
    Solution NoSQL choisie.
    Architecture mise en place.
    Résultats (performances, coûts).

    \note{
        Présenter le cas en quatre temps : besoin métier, données à stocker,
        architecture retenue, puis bénéfices mesurables.
        Faire identifier la clé d'accès et vérifier que les opérations sont
        majoritairement des lectures ou écritures ciblées.
        Ne pas présenter Redis comme un remplacement universel d'une base durable :
        discuter la persistance, la réplication et le coût de la mémoire.
    }

\end{frame}
